\section{Bidirectional signed forcing at the original roof}
\label{sec:v95-bidirectional}
A positive entrance is a creation term in the backward comparison but
a loss term in the forward comparison. We retain both comparisons,
including all opposing physical and artificial cuts, and take their
minimum only after projection. This does not reverse or reindex the
billiard record.

Write $R=R_c^\varepsilon$ and, for a fixed $\rho>0$, put
\begin{align}
 \mathsf H_\rho^+&=\mathsf G_c^{\varepsilon,+,\ge\rho},
 &\mathsf H_\rho^-&=\mathsf G_c^{\varepsilon,-,\ge\rho},\notag\\
 \mathsf K_\rho^-&=DR-\mathsf H_\rho^+,
 &\mathsf K_\rho^+&=-DR-\mathsf H_\rho^-.
 \label{eq:v95-two-measures}
\end{align}
Superscripts on $\mathsf K$ indicate the side of the target; those
on $\mathsf H$ indicate the sign of the density derivative. Thus
\[
 \mathsf K_\rho^+=\mathsf D_c^\varepsilon
 +\mathsf G_c^{\varepsilon,-,<\rho}
 -\mathsf I_c^\varepsilon-\mathsf G_c^{\varepsilon,+}.
\]
In this forward measure a critical entrance is negative. An upper
prefix cut is positive and is not discarded.

Let $A_r^-R=r^{-1}\int_{t-r}^tR$ and
$A_r^+R=r^{-1}\int_t^{t+r}R$. For finite signed measures define
\[
 \mathcal T_r^-\mu(t)=\int_{(t-r,t]}\frac{u-t+r}{r}\,d\mu(u),
 \qquad
 \mathcal T_r^+\mu(t)=\int_{[t,t+r)}\frac{t+r-u}{r}\,d\mu(u).
\]
At Lebesgue roofs of $R$ and away from atoms, set
\begin{align}
 S_{r,\rho}^-&=[\mathcal T_r^-\mathsf K_\rho^-]_+,
 &S_{r,\rho}^+&=[\mathcal T_r^+\mathsf K_\rho^+]_+,\notag\\
 J_{c;r,\rho}^\varepsilon&=\min\{S_{r,\rho}^-,S_{r,\rho}^+\}.
 \label{eq:v95-joint-forcing}
\end{align}
Define $S^{-,\sharp}$ and $S^{+,\sharp}$ by taking the essential
supremum of $[\mathsf K^-((s,t])]_+$ over $t-r<s<t$ and
$[\mathsf K^+([t,s))]_+$ over $t<s<t+r$, respectively; put
$J^\sharp=\min\{S^{-,\sharp},S^{+,\sharp}\}$. Finite-count BV
representatives give Borel versions. No common Lebesgue-point set
for an uncountable collection of parameters is asserted.

\begin{theorem}[A smaller two-sided pointwise budget]
\label{thm:v95-bidirectional}
For every previously fixed $r,\rho>0$, at finite count and almost
every original roof,
\begin{align}
 0\le J_{c;r,\rho}^\varepsilon
 &\le J_{c;r,\rho}^{\varepsilon,\sharp}\le R_c^\varepsilon,
 \label{eq:v95-joint-domination}\\
 R_c^\varepsilon(t)
 &\le \max\{A_r^-R_c^\varepsilon(t),A_r^+R_c^\varepsilon(t)\}
       +J_{c;r,\rho}^\varepsilon(t)
       +C E_\varepsilon\int_{t-r}^{t+r}v^\rho(u)\,du.
 \label{eq:v95-joint-pointwise}
\end{align}
The joint forcing is bounded by the backward signed forcing of
module 225. If its sharp version plus the last error is at most
$R_c^\varepsilon(t)-\lambda$, then the source is at least $\lambda$
almost everywhere on at least one of the two full prescribed adjacent
intervals. All sources, losses and positive parts are formed at finite
count, before any uniform height is taken.
\end{theorem}
\begin{proof}
Stieltjes integration of the two identities in
\eqref{eq:v95-two-measures} gives
\begin{align*}
 R-A_r^-R&=\mathcal T_r^-\mathsf K_\rho^-
                         +\mathcal T_r^-\mathsf H_\rho^+,\\
 R-A_r^+R&=\mathcal T_r^+\mathsf K_\rho^+
                         +\mathcal T_r^+\mathsf H_\rho^-.
\end{align*}
The signs can also be checked by integrating the fundamental theorem
between the target and the earlier or later roof. Each last term is
nonnegative. Each signed interval increment, after subtraction of
its $\mathsf H$ term, is at most $R(t)$ because the other endpoint
source is nonnegative. This proves the domination for both sharp
forcings. Their ramped expressions are averages of the corresponding
interval increments, so $S^\pm\le S^{\pm,\sharp}$. Taking minima
proves \eqref{eq:v95-joint-domination}.

The derivative bound, with its physical angular density retained,
gives
\[
 \mathsf H_\rho^+((t-r,t])+\mathsf H_\rho^-([t,t+r))
 \le C E_\varepsilon\int_{t-r}^{t+r}v^\rho(u)\,du.
\]
This follows from multiplicity-one coarea summation, not from a
component count. Either ramped $\mathsf H$ term is at most this
common error. Take the smaller of the two upper comparisons and use
$\min\{x+s,y+t\}\le\max\{x,y\}+\min\{s,t\}$.
This proves \eqref{eq:v95-joint-pointwise}.
The two interval identities likewise give lower bounds for $R(s)$
on the backward and forward interval. Choose the side attaining the
smaller sharp forcing to obtain the asserted persistence statement.

All physical cuts and all artificial prefix cuts are already present
in $DR$. The positive parts in \eqref{eq:v95-joint-forcing} are
therefore taken after their signed projection, not sheet by sheet.
In particular a critical entrance trace alone cannot be positive
forcing in both directions. Forward terminations and negative critical
growth can still make the forward forcing large.
\end{proof}

The minimum of two projected forcings is not the integral of the
minimum of two sheetwise forcings. No such interchange is used.
The packet $M$ of Section~\ref{sec:v95-packets} is a separate positive
majorant with its own proved error bound; it is not identified with
$J$ or with a Jordan part of $DR$.

\subsection{The original scalar endpoint and legal rates}
For the full bad source take $c_\sigma=b_\sigma$. Define
\[
 \mathfrak J_{\varepsilon,M}(r)
   =\lim_{\rho\downarrow0}\mathcal H_M(J_{r,\rho}^\varepsilon),
 \qquad
 \mathfrak J_{\varepsilon,M}^{\sharp}(r)
   =\lim_{\rho\downarrow0}\mathcal H_M(J_{r,\rho}^{\varepsilon,\sharp}).
\]
Both limits are monotone extended nonnegative limits: lowering $\rho$
increases the removed positive measures $\mathsf H_\rho^\pm$.

\begin{theorem}[Bidirectional criterion for the unchanged pointwise law]
\label{thm:v95-endpoint}
For all sufficiently large fixed $B$ and every fixed $r>0$,
\begin{equation}\label{eq:v95-endpoint}
 \mathcal E_M\le C_MB^{-1/192}
 +\mathfrak J_{\varepsilon(B),M}(r)
 +C E_{\varepsilon(B)}r.
\end{equation}
For any prescribed $r(B)>0$ with $E_{\varepsilon(B)}r(B)\to0$,
the original pointwise assertion is equivalent to either
$\mathfrak J_{\varepsilon(B),M}(r(B))\to0$ or
$\mathfrak J_{\varepsilon(B),M}^{\sharp}(r(B))\to0$, for every fixed $M$.
\end{theorem}
\begin{proof}
For a nonnegative source both one-sided averages are at most
$2A_rb^\varepsilon$. Fix $B,r,\rho$ and the auxiliary bands;
apply the original fixed-width bad-source estimate, the cutoff
volume estimate on length $2r$, and the scalar positive remainder.
Only after the collision limsup and allowed band enlargement do we
let $\rho$ decrease. Since
$\varepsilon(B)^{1/16}=O(B^{-1/192})$, the result is
\eqref{eq:v95-endpoint}. This proves sufficiency of either criterion.
Conversely the original pointwise law and its positive remainder give
$\mathcal H_M(b^{\varepsilon(B)})\le C_MB^{-1/192}$.
Domination \eqref{eq:v95-joint-domination} proves necessity for every
prescribed radius, without exchanging limits.
\end{proof}

The explicit new errors admit a common legal schedule. With
$r(B)=c\varepsilon(B)^N$, $q(B)=\sqrt{r(B)}$ and $N>1$,
\[
 E_{\varepsilon(B)}r(B)\to0,\qquad
 q(B)^\alpha+E_{\varepsilon(B)}q(B)^{\alpha+1}\to0.
\]
The same schedule pays the grazing-ended short-packet error when
$\delta=r(B)$ at each previously fixed positive entrance gradient.
These rates do not prove that the remaining two-nongrazing packet
charge or the zero-gradient target concentration vanishes. Nor do they
justify an infinite radial sum inside a collision limsup. The theorem
is an improved pointwise-controlling criterion, not its verification.
The full scalar law and unrestricted point-conditioned consequences
remain the governing objectives, with their original arithmetic
restrictions intact.
